\subsection{Produktkrav}
\label{sec:produktkrav}
\subsubsection{Kravgrunnlag og prioritering}
Kravspesifikasjonen er et arbeidsutkast. Oppgitte rammer fra oppgavebeskrivelsen kontrolleres mot gjeldende avtale med veileder før designet låses. Egne ytelsesmål holdes adskilt fra obligatoriske rammer.

\subsubsection{Foreløpig kravspesifikasjon}
\begin{description}
    \item[R1--R4: Oppgitte rammer.] Ytre mål skal ikke overstige $140 \times 85 \times 70$ mm. Utdelt stålsvinghjul med diameter 60 mm og tykkelse 6 mm benyttes, med oppgitt svinghjulaksel med diameter 6 mm og lengde 65 mm. Opptrekk utføres med 20 cm tråd. Verifiseres ved måling og inspeksjon mot oppgavegrunnlaget.
    \item[R5: Funksjon.] Bilen skal kjøre selvstendig etter opptrekk. Verifiseres ved funksjonstest.
    \item[R6: Hovedmål.] Størst mulig gjennomsnittlig kjøredistanse ved standardisert opptrekk. Vurderes med minst 10 kjøringer per sammenlignet konfigurasjon.
    \item[R7: Repeterbarhet.] Spredning i kjøredistanse dokumenteres med standardavvik og eventuelt variasjonskoeffisient. Akseptgrense fastsettes etter referansetestene, før senere sammenligning.
    \item[R8: Retningsstabilitet.] Sideavvik måles med en fast prosedyre. Akseptgrense fastsettes etter referansetestene.
    \item[R9: Pålitelighet.] Eget foreløpig mål: minst 9 av 10 opptrekk gir vellykket kjøring. Kriteriet for vellykket kjøring defineres før testserien.
    \item[R10--R12: Justerbarhet.] Drivverket bør kunne demonteres og justeres. Utveksling og hjul bør kunne endres uten fullstendig redesign. Verifiseres ved inspeksjon og praktisk utprøving.
    \item[R13: Produksjon.] Bilen skal kunne fremstilles med tilgjengelige produksjonsmetoder. Verifiseres mot produksjonsunderlag og ferdig prototype.
\end{description}
\subsubsection{Verifikasjon og revisjon av krav}
% Registrer krav-ID, endelig krav, kilde, prioritet, testreferanse og status.
Det settes ikke et vilkårlig minstekrav til kjøredistanse før referansetester foreligger. Endringer i krav begrunnes og dateres.
